% NS-43 analytic working note. No numerical computation or certificate.
% Scope: additional complete-space inputs for CCM step (b).

\subsection*{CCM selection requires a relative estimate, not another absolute residual}

\paragraph{Status.}
The fixed-window simple ground and its real-zero transform do not identify
Riemann's $\Xi$. The known source limit and complete residual are useful
prerequisites. The new input would be a cofinal \emph{complete spectral
selection estimate}, with a rate compatible with the transform topology.
The propositions below are conditional elementary functional analysis;
none establishes the hypotheses for the arithmetic Weil family.

Put $H=L^2(\mathbb R)$, $I_a=(-a,a)$ and let $J_a$ denote zero extension.
Let $A_a$ be the complete lower-bounded self-adjoint even-sector operator
on $L^2(I_a)$, with compact resolvent. If the global ground is used,
its evenness and ordering below the odd sector must also be supplied.
In the following selection lemmas assume the selected ground is simple,
with unit vector $v_a$, eigenvalue $\mu_a$, and complete next-level
separation $g_a>0$:
\[
 q_a[f]-\mu_a\|f\|^2\ge g_a\|f\|^2
 \quad(f\in\mathcal D(q_a),\ f\perp v_a).
\]
This is an ordinary complete spectral separation, not an LDL pivot,
finite-compression gap, or a source-complement estimate inferred from
separate diagonal blocks.

Let $h$ be the explicit positive Gaussian arithmetic radical of
\textup{eq:v150-weight}, and $\phi=h/\|h\|_2$. Use exactly the fixed
boundary profile of \textup{prop:v135-growing-radical}: fix one smooth
$\vartheta$ with $0\le\vartheta\le1$, $\vartheta(t)=1$ for $t\le-1$,
and $\vartheta(t)=0$ for $t\ge-1/2$. For $a\ge2$ set
$\chi_a(x)=\vartheta(|x|-a)$. This is smooth (it is constant near zero),
even, supported in $I_a$, equal to one on $[-a+1,a-1]$, and has
uniformly bounded derivatives of every fixed order. Define
\[
 p_a=\chi_a\phi/\|\chi_a\phi\|_2,\qquad
 \alpha_a=q_a[p_a],\qquad E_a=\alpha_a-\mu_a\ge0.
\]
For large $a$, these vectors lie in the complete operator domain.
The double-exponential tail from v1.35 implies, for every fixed $R\ge0$,
\[
 \int_{\mathbb R}e^{R|x|}|J_ap_a(x)-\phi(x)|\,dx\longrightarrow0.
\]
The vector $p_a$ is the normalized $j=0$ column of the growing block.
Thus \textup{prop:v135-growing-radical}, specifically
\textup{eq:v135-block-residual}, gives the complete estimate
\[
 \|A_ap_a\|\le\varepsilon_a:=Ca\exp(-e^a/100),\qquad
 |\alpha_a|\le\varepsilon_a,
\]
with constants for this fixed source and profile. No residual estimate
is asserted for arbitrary cutoffs whose derivatives may grow with $a$.
The cutoff $p_a$ here is explicit and is not silently identified with
the different repaired prolate source $u_a$.

\begin{proposition}[A complete relative selection estimate]
\label{prop:ns43-ccm-rayleigh-selection}
Choose the phase of $v_a$ so that $\langle p_a,v_a\rangle\ge0$.
If $E_a/g_a\to0$, then $\|J_av_a-\phi\|_2\to0$.
Define
\[
 C_R(a)=\begin{cases}
 \sqrt{(e^{2Ra}-1)/(2\pi R)},&R>0,\\
 \sqrt{a/\pi},&R=0.
 \end{cases}
\]
If, in addition, for every fixed $R>0$,
\[
 C_R(a)\sqrt{E_a/g_a}\longrightarrow0,
 \tag{S1}
\]
then the unitary Fourier transforms of $J_av_a$ converge locally
uniformly on $\mathbb C$ to $\widehat\phi$.
\end{proposition}
\begin{proof}
Write $p_a=\kappa_a v_a+e_a$, where $e_a\perp v_a$ and
$\kappa_a\ge0$. The spectral projection preserves the form domain,
and the cross term is exactly zero because $v_a$ is an eigenvector.
Consequently
\[
 E_a=q_a[e_a]-\mu_a\|e_a\|^2\ge g_a\|e_a\|^2,
 \quad \kappa_a^2=1-\|e_a\|^2,
 \quad \|p_a-v_a\|^2=2(1-\kappa_a)\le2E_a/g_a.
\]
This proves the ordinary-norm conclusion. For $|\Im z|\le R$,
Cauchy--Schwarz gives
\[
 |\widehat{J_a(v_a-p_a)}(z)|
 \le (2\pi)^{-1/2}\left(\int_{-a}^a e^{2R|x|}\,dx\right)^{1/2}
       \|v_a-p_a\|
 \le C_R(a)\sqrt{2E_a/g_a}.
\]
The weighted cutoff convergence above controls
$\widehat{J_ap_a}-\widehat\phi$ on the same strip. Combining the
bounds proves the assertion. No compactness hypothesis is hidden:
(S1) supplies a direct quantitative topology upgrade.
\end{proof}

\begin{proposition}[A residual-to-separator version]
\label{prop:ns43-ccm-residual-selection}
Let $p_a\in\mathcal D(A_a)$ be unit, put
$\alpha_a=\langle A_ap_a,p_a\rangle$ and
$\rho_a=\|(A_a-\alpha_a)p_a\|$. Suppose an independent complete
second-eigenvalue lower bound $\beta_a>\alpha_a$ is established,
with eigenvalues counted with multiplicity. Then the ground is simple and
\[
 \|v_a-p_a\|\le\frac{\sqrt2\,\rho_a}{\beta_a-\alpha_a}.
\]
Thus the cofinal condition
\[
 C_R(a)\frac{\rho_a}{\beta_a-\alpha_a}\longrightarrow0
 \quad\text{for every fixed }R>0 \tag{S1r}
\]
selects the entire transform of the explicit cutoff radical.
\end{proposition}
\begin{proof}
The variational upper bound is $\mu_a\le\alpha_a<\beta_a$, so
exactly one eigenvalue lies below $\beta_a$ and it is simple.
On its orthogonal complement the spectrum of $A_a-\alpha_a$ lies
above $\beta_a-\alpha_a$. Spectral expansion gives
$\rho_a^2\ge(\beta_a-\alpha_a)^2\|e_a\|^2$.
The preceding phase/norm identity and Fourier estimate complete the proof.
\end{proof}
This formulation does not require a lower ground-energy estimate matching
$\alpha_a$ to high relative precision. The missing arithmetic input is
the complete cofinal separator at the correct relative scale. A finite
matrix's second eigenvalue supplies no such lower bound without its full
tail and coupling comparison. An even-sector separator still needs the
separate odd ordering if the real-zero theorem requires the global ground.

\paragraph{A weaker alternative to the evaluator rate.}
Ordinary selection $E_a/g_a\to0$ also gives locally uniform entire
convergence if the selected grounds satisfy the independently proved
weighted tightness condition
\[
 \lim_{K\to\infty}\sup_a\int_{|x|>K}e^{R|x|}|J_av_a(x)|\,dx=0
 \quad\text{for every }R>0. \tag{S2}
\]
Indeed the integral on $[-K,K]$ converges by ordinary $L^2$ convergence
and Cauchy--Schwarz, while (S2) and the known tail of $\phi$ control
the exterior. A sufficient explicit condition for (S2) is
$\sup_a\int e^{2(R+1)|x|}|J_av_a|^2<\infty$ for each $R$.
No such ground-tail statement follows from the current residual bound.

\begin{proposition}[Complete evaluator-dual alternative]
\label{prop:ns43-ccm-evaluator-selection}
In the preceding setting suppose $E_a/g_a\to0$. Suppose further that,
for every fixed $R>0$, a complete complement estimate is established:
\[
 \sup_{|z|\le R}|\widehat{J_ae}(z)|^2
 \le D_{a,R}\bigl(q_a[e]-\mu_a\|e\|^2\bigr)
 \quad(e\in\mathcal D(q_a),\ e\perp v_a),
 \qquad D_{a,R}E_a\to0. \tag{S3}
\]
Then the same locally uniform entire convergence follows.
\end{proposition}
\begin{proof}
For the orthogonal error above,
$\widehat p_a=\kappa_a\widehat v_a+\widehat e_a$ and
$\kappa_a\to1$. Thus, on the disk,
\[
 |\widehat v_a-\widehat p_a|
 \le \kappa_a^{-1}\sqrt{D_{a,R}E_a}
       +(\kappa_a^{-1}-1)|\widehat p_a|\to0.
\]
The explicit cutoff transforms are uniformly bounded on each disk
and converge to $\widehat\phi$. This proves the assertion.
\end{proof}
The shift by $\mu_a$ and the complete complement in (S3) are part of
its hypotheses. The fixed-window v1.43 evaluator comparison for one
real zero is not a cofinal bound of this kind, nor a bound on every
complex compact set. Its scalar resolution floor is not a no-go
result for (S1)--(S3).

\begin{proposition}[A source-relative block condition sufficient for selection]
\label{prop:ns43-ccm-source-separation}
Let $p_a\in\mathcal D(A_a)$ be unit, $P_a=|p_a\rangle\langle p_a|$,
$Q_a=I-P_a$, $\alpha_a=\langle A_ap_a,p_a\rangle$, and
$r_a=Q_aA_ap_a$. Suppose the \emph{complete compressed form} satisfies
\[
 q_a[f]\ge(\alpha_a+\delta_a)\|f\|^2
 \quad(f\in\mathcal D(q_a)\cap p_a^\perp),\qquad \delta_a>0.
 \tag{S4}
\]
Then $A_a$ has exactly one eigenvalue below $\alpha_a+\delta_a$,
it is simple, and, for its phase-aligned ground $v_a$,
\[
 \alpha_a-\|r_a\|^2/\delta_a\le\mu_a\le\alpha_a,
 \qquad \|v_a-p_a\|\le\sqrt2\,\|r_a\|/\delta_a.
\]
For the explicit cutoff source, the cofinal rate
$C_R(a)\|r_a\|/\delta_a\to0$ for every $R>0$ therefore implies
entire-transform selection.
\end{proposition}
\begin{proof}
The restriction to $p_a^\perp$ is a closed lower-bounded form,
because subtracting its component along the fixed operator-domain
vector preserves the form domain continuously. The min--max principle
places the second eigenvalue at least $\alpha_a+\delta_a$; the trial
$p_a$ puts the lowest at most $\alpha_a$. For $f=tp_a+y$, $y\perp p_a$,
\[
 q_a[f]-\alpha_a\|f\|^2
 =q_a[y]-\alpha_a\|y\|^2+2\Re(t\langle A_ap_a,y\rangle)
 \ge\delta_a\|y\|^2-2|t|\|r_a\|\|y\|
 \ge-\|r_a\|^2|t|^2/\delta_a.
\]
This handles the full cross term and proves the lower edge bound.
Writing $v_a=tp_a+y$, the weak projected eigen-equation is
$(Q_aA_aQ_a-\mu_a)y=-tr_a$. Its inverse has norm at most
$1/(\alpha_a+\delta_a-\mu_a)\le1/\delta_a$, so
$\|y\|\le|t|\|r_a\|/\delta_a$. Phase alignment then gives
$\|v_a-p_a\|^2=2(1-|t|)\le2\|y\|^2$.
\end{proof}
Condition (S4) is a new, source-specific signed spectral estimate,
not an algebraic reformulation of the bounded-floor target. It asks for information beyond a floor: it separates one named source
from its whole complement and controls the relative error. No reverse
implication from RH or from a uniform floor to (S4) is claimed. In this
arithmetic family, with the known vanishing source residual, a successful
cofinal condition of this strength gives a uniform floor in the sector
where it is imposed. To use v1.36 one must impose it on the full space
or separately control the odd sector by a uniform floor; those additional
hypotheses would imply RH. Even-sector selection does not silently supply
the odd estimate. It is not an unconditional shortcut.

The v1.35 even reflected near-radical block has dimension $J_a+1$,
so eventually it contains a unit vector orthogonal to $p_a$. Its complete
residual is at most $\varepsilon_a$, hence (S4) forces
$\delta_a\le\varepsilon_a-\alpha_a\le\varepsilon_a+|\alpha_a|$.
The explicit cutoff source also has $|\alpha_a|=O(\varepsilon_a)$.
This rules out a uniform positive $\delta_a$ or a polynomial lower scale.
A source-selecting separation must shrink and must still dominate the
source error in the relative sense above. The existence of a growing
near-zero block does not logically exclude such a hierarchy. Neither
that hierarchy nor its required rate is in the archive.

\paragraph{Identification of the limit.}
By \textup{eq:v150-mellin}, with
$\Xi(z)=\xi(1/2+iz)$ and the even functional equation,
\[
 \widehat\phi(z)=\frac{2}{\sqrt{6\pi}\,\|h\|_2}\,\Xi(z).
\]
In particular $\widehat\phi(0)>0$. Once either selection criterion is
proved, $\widehat v_a(0)\ne0$ eventually and
\[
 \frac{\widehat v_a(z)}{\widehat v_a(0)}
 \longrightarrow\frac{\Xi(z)}{\Xi(0)}
 \quad\text{locally uniformly on }\mathbb C.
\]
The scalar is fixed by the stated $L^2$ normalization and then by the
value at zero; phase/scale and translations must not be left free.
If the real-zero-ground-transform theorem's complete hypotheses hold
cofinally as well, this convergence rules out a nonreal zero of $\Xi$:
a small closed disk about such a zero, disjoint from the real axis and
with zero-free boundary, would by uniform convergence and the argument
principle eventually contain the same positive number of zeros of
$\widehat v_a$, a contradiction. This is a conditional implication,
not a new RH proof. Cofinal simplicity/evenness and applicability of
that theorem are themselves additional to the single $\lambda=4$ result.

\subsection*{An exact abstract nonselection model}

The next construction tests precisely the proposed \emph{abstract}
premises: complete compact-resolvent operators on increasing intervals;
simple even nonnegative ground functions; entire ground transforms with only real
zeros; monotone positive ground energies tending to zero; a fixed
candidate with vanishing full residual; dense fixed tests with vanishing
residual; ordinary strong-resolvent collapse; and even uniformly bounded
support of the ground functions. It does \emph{not} impose the
support-consistent arithmetic Weil form, the prime identities, the
CvS distribution representation, or a positivity-improving semigroup.
Its ground functions are nonnegative, not strictly positive on all of
the increasing interval. It cannot disprove an arithmetic
selection theorem using those additional structures.

\begin{proposition}[Abstract nonselection despite fast absolute residuals]
\label{prop:ns43-ccm-nonselection}
There is a family satisfying all the abstract properties just listed,
including uniformly tight ground supports and the explicit Gaussian
cutoff candidate, whose phase-fixed normalized ground transforms have
two different subsequential limits. The absolute residual rate can be
made arbitrarily fast by a positive scalar rescaling.
\end{proposition}
\begin{proof}
Let $b=1_{[-1/4,1/4]}$ and let $g_1=b*b*b*b/\|b*b*b*b\|_2$.
Then $g_1$ is nonnegative and even, lies in $H^2(\mathbb R)$, and is
supported in $[-1,1]$. Put $g_2(x)=2^{-1/2}g_1(x/2)$, supported in
$[-2,2]$. Both are unit vectors, distinct and not scalar multiples.
Their entire transforms are nonzero constants times
$\operatorname{sinc}(z/4)^4$ and $\operatorname{sinc}(z/2)^4$,
respectively, where $\operatorname{sinc}z=\sin z/z$. Every zero is real.
Their normalized transforms have distinct zero sets.

For integers $n\ge3$, let $D_n=-d^2/dx^2$ be the Dirichlet Laplacian
on $H_n=L^2(-n,n)$, with domain $H^2(-n,n)\cap H_0^1(-n,n)$.
Its unit first eigenvector is
$h_n(x)=n^{-1/2}\cos(\pi x/(2n))$.
Let $j(n)=1$ for even $n$ and $2$ for odd $n$. Define
\[
 w_n=\frac{h_n-g_{j(n)}}{\|h_n-g_{j(n)}\|},\qquad
 U_n=I-2|w_n\rangle\langle w_n|,\qquad
 A_n=n^{-1}U_n(I+D_n)U_n.
\]
The denominator is nonzero; its square is
$2-2\langle h_n,g_{j(n)}\rangle\to2$.
Thus $U_n$ is an even-parity-preserving unitary involution with
$U_nh_n=g_{j(n)}$. Because $w_n\in\mathcal D(D_n)$, it preserves
that domain. The complete $A_n$ is positive, self-adjoint with compact
resolvent, and its simple even ground is exactly $g_{j(n)}$ with
\[
 \mu_n=\frac1n\left(1+\frac{\pi^2}{4n^2}\right)\downarrow0,
 \qquad \lambda_{2,n}-\mu_n=\frac{3\pi^2}{4n^3}>0.
\]
All ground transforms have only real zeros, but the normalized ground
functions and their normalized entire transforms alternate between two
different limits. Positivity and a fixed positive phase do not repair this.

For each fixed $f\in C_c^\infty(\mathbb R)$ and all sufficiently large
$n$, $f\in\mathcal D(D_n)$. Moreover
$\|(I+D_n)w_n\|$ is uniformly bounded: the $g_j$ have fixed finite
$H^2$ norms, while $\|D_nh_n\|=\pi^2/(4n^2)$ and the denominators
are bounded away from zero. Consequently
\[
 \|A_nf\|\le\frac1n\left(\|f-f''\|
        +2\|f\|\,\|(I+D_n)w_n\|\right)\longrightarrow0.
\]
The same estimate holds for the fixed candidate $g_1$. It also holds
uniformly for the explicit Gaussian cutoff candidate $p_n$ used above:
its $H^2$ norms are uniformly bounded, its cutoff normalization stays
away from zero, and $\|p_n\|=1$. Thus $\|A_np_n\|=O(1/n)$ while
$p_n\to\phi$ in every exponentially weighted $L^1$ norm. Even the
known explicit Gaussian candidate and its transform identity do not
select the alternating grounds of this abstract operator family.
Let $\widetilde A_n=A_n\oplus0$ on $L^2(\mathbb R)$ by zero extension.
For every nonreal $z$ and every such fixed $f$,
\[
 \|(\widetilde A_n-z)^{-1}f+z^{-1}f\|
 \le\frac{\|\widetilde A_nf\|}{|z||\Im z|}\to0.
\]
Density of $C_c^\infty$ and the uniform self-adjoint resolvent bound
extend this to all of $H$. Thus the ordinary strong-resolvent limit
is zero. This model meets every listed abstract premise, including
weighted tightness of the grounds, yet supplies no unique selection.
Its ground multiplicity is one at every $n$; the failure is not caused
by a degenerate eigenspace at a fixed window.
The scalar $1/n$ can be replaced throughout by any positive sequence
$\theta_n\downarrow0$. The eigenvalues and residuals are rescaled,
the ground functions are unchanged, and the residual estimate becomes
$O(\theta_n)$. Thus even an arbitrarily prescribed rapid absolute
residual rate cannot replace the missing relative estimate. This
rescaling does not restore arithmetic support consistency.
\end{proof}


\paragraph{Ordinary norm is also the wrong transform topology by itself.}
Even if a selection estimate gives $f_a\to\phi$ in $L^2$, locally
uniform entire convergence still needs a rate or a tail condition.
For any fixed $R>0$, add to the positive cutoff source a positive even
unit bump supported near $\pm(a-1)$ with coefficient $e^{-Ra/2}$,
and renormalize. The resulting compactly supported unit functions
converge to $\phi$ in $L^2$, while their transforms at $iR$ grow like
a positive constant times $e^{Ra/2}$. This is a separate topology
counterexample; it does not claim the real-zero property or the
arithmetic ground equation.

\subsection*{A restriction on renormalized profile selection}

A possible alternative is to renormalize the shrinking spectrum and
prove a nonzero profile form with a unique selected direction.
The ordinary zero-resolvent limit cannot provide this profile.
Here is a precise obstruction to choosing too coarse a scale.
\begin{proposition}[A coarse rescaled liminf cannot select one radical]
\label{prop:ns43-ccm-coarse-profile}
Let $s_a>0$ and suppose
\[
 \varepsilon_a/s_a\to0,\quad |\alpha_a|/s_a\to0,
 \qquad \varepsilon_a=Ca e^{-e^a/100},
\]
where each fixed finite translate combination may have its own constant
$C$, as in v1.35--v1.36. Suppose a nonnegative lower-semicontinuous
quadratic form $\mathcal E$ on $H$ satisfies the proposed strong
liminf inequality
\[
 \mathcal E(f)\le\liminf_a
 \frac{q_a[f_a]-\alpha_a\|f_a\|^2}{s_a}
 \quad\text{whenever }J_af_a\to f.
 \tag{S5}
\]
Then $\mathcal E\equiv0$ on $H$.
\end{proposition}
\begin{proof}
For every fixed $f$ in the dense translated-radical span, insert its
smooth cutoff $f_a$. The complete residual implies the right side is
zero, so $\mathcal E(f)=0$. Lower semicontinuity and density then imply
$\mathcal E\equiv0$ on $H$. Therefore no such coarse rescaled limit
can have one-dimensional kernel $\mathbb C\phi$.
\end{proof}
This is a scoped implication about (S5), not a prohibition on finer
scales, different compactness mechanisms, or arithmetic ground selection.

\paragraph{Recommended node readings.}
\begin{itemize}
\item CCM step (b): retain the map status \textbf{live} (the sole gold
node), as explicitly requested by the user; its mathematical target remains
\textbf{open}. A finite-window real-zero ground and
one fixed-window zero enclosure are retained; cofinal function selection
and its entire-transform topology are missing.
\item Gaussian source identity and residual: \textbf{proved prerequisite}.
The identity selects an explicit candidate, not the operator's ground.
\item Absolute-residual / ordinary-resolvent inference to a unique
Gaussian ground: \textbf{closed as an inference from the listed abstract
premises}, by the exact countermodel. Do not label the arithmetic CCM
route closed.
\item Complete relative spectral/evaluator selection, (S1) or (S3), or
the stronger source block condition (S4): \textbf{blocked/open input}.
No cofinal separation, Rayleigh-excess comparison, or ground-tail rate
is presently available.
\item Coarse renormalized single-profile limit satisfying (S5):
\textbf{closed only under its explicit scale and liminf hypotheses}.
Finer renormalized selection remains open.
\end{itemize}
Finally, in the actual support-consistent Weil family, a uniform lower
bound on the \emph{complete} ground energy is already the bounded-floor
target of v1.36, and a ground-energy limit zero is equivalent to positivity
via monotonicity and the known source upper bound. Those must not be
rebranded as weak new selection assumptions. Neither conclusion alone
has been shown here to identify the Gaussian ground profile. The
abstract countermodel establishes no arithmetic nonimplication.
